\hypertarget{r19.-removing-the-class-number-one-hypothesis-mean-value-formulas-for-arbitrary-module-lattices-over-a-maximal-order}{%
\subsection*{R19. Removing the class-number-one hypothesis: mean value formulas for arbitrary module lattices over a maximal order}\label{r19.-removing-the-class-number-one-hypothesis-mean-value-formulas-for-arbitrary-module-lattices-over-a-maximal-order}}

\begin{quote}
\textbf{STATUS (written in one pass under a 24 minute limit, NOT refereed, no numerics).}

\begin{itemize}
\tightlist
\item
  \textbf{Part A (general rank-one formula): PROVED {[}P{]}} from (L1)\(_m\) used as a black box (statement only, for Borel \(f\ge0\)), plus uniqueness of invariant Radon measures on a homogeneous space, Fatou\textquotesingle s lemma and Riemann sums. No new domination estimate is needed (averaging trick, A.4).
\item
  \textbf{Part B (\(j\)-frames for \(M L_0\), any Steinitz classes): PROVED {[}P{]}} from Part A in ranks \(k,\dots,k-j+1\) (hence from (L1)\(_m\), \(k-j+1\le m\le k\)), for \(\mathcal O\) hereditary (in particular maximal).
\item
  \textbf{Part C (Theorem E\textquotesingle): PROVED {[}P{]} at the level of the mean value formula (MVF)} (Lemma 1.2\textquotesingle). All of R14 \texttt{proof.md} (§0-§7) was then read in this session \emph{for uses of freeness / of the coordinates \(\Lambda_W=\mathcal O^r\)} (result in C.4: none beyond §0.1, Lemma 0.5, Lemma 1.1, Lemma 1.2 and the notation "\(M\mathcal O^{r\times k}\gamma\)"). The analytic correctness of R14 §2-§5 was NOT re-verified here; Theorem E\textquotesingle{} inherits it.
\item
  \textbf{No {[}GAP{]} is known to me in Parts A, B, C.} Cited and unverified: (L1) (Gargava, status as in R17), (Her) (Reiner, from memory), (U), Weil\textquotesingle s quotient formula, (L2), (L3).
\end{itemize}
\end{quote}

Tags: \textbf{{[}P{]}} proved here; \textbf{{[}L{]}} cited (stated whether read in this session or from memory); \textbf{{[}GAP{]}}.
Sources read in this session: \texttt{../R18\_integration/report.md} (all), \texttt{../R15\_rogers/proof.md} (all), \texttt{../R14\_leech\_tensor/proof.md} (all, §2-§7 only scanned for the question of C.4), \texttt{../R17\_gargava\_check/notes.md} (all), \texttt{../R17\_gargava\_check/referee.md} (first 37 lines: verdict CORRECT\_MODULO\_MINOR\_FIXES, confirms (R1)-(R4) and Lemma 18). NOT read: Gargava\textquotesingle s paper itself, Weil\textquotesingle s notes, Reiner\textquotesingle s book, Klartag\textquotesingle s paper, the R5/R6/R10 files.

\hypertarget{setting}{%
\subsubsection*{0. Setting}\label{setting}}

\(\Delta\) a \(\mathbb Q\)-division algebra, \(D:=\dim_{\mathbb Q}\Delta\), \(\Delta_{\mathbb R}=\prod_v\mathbb K_v\), \(\mathbb K_v\in\{\mathbb R,\mathbb C,\mathbb H\}\); \(\mathcal O\subset\Delta\) a \textbf{maximal} order. Rows \(\Delta_{\mathbb R}^{1\times m}\): scalars act on the left, matrices on the right. \(\|\gamma\|:=|\det_{\mathbb R}(y\mapsto y\gamma)|\) on rows (\(=\) the same on columns, {[}R15, Lemma 0.2{]}). \(\mathsf G_m:=\{\gamma\in GL_m(\Delta_{\mathbb R}):\|\gamma\|=1\}\) (normal in \(GL_m(\Delta_{\mathbb R})\): kernel of a character), \(\Gamma_m:=GL_m(\mathcal O)\subset\mathsf G_m\). Lebesgue measure on \(\Delta_{\mathbb R}^{a\times b}\) normalised by \(\mathrm{covol}(\mathcal O^{a\times b})=1\) (product-compatible). \(\Omega_m:=\prod_v(\mathbb K_v^{1\times m}\setminus0)\).

\textbf{(Her) {[}L, from memory: Reiner, Maximal Orders, Thm 21.4 + Cor. 2.44/Thm 2.44{]}} A maximal order is left and right hereditary; hence every finitely generated torsion-free left (right) \(\mathcal O\)-module is projective, and every surjection of an \(\mathcal O\)-module onto such a module splits. \emph{Only this consequence of maximality is used below.}

\textbf{(L1)\(_m\), \(m\ge2\) {[}L; Gargava arXiv:2107.04844v3 Thm 16 + Prop. 14, as read and repaired in R17 \texttt{notes.md}; read there, not in this session{]}} \(\Gamma_m\) is a lattice in \(\mathsf G_m\), and for Borel \(f\ge0\) (extension from \(C_c\): {[}R15, Lemma 0.5{]}) \(\int_{\Gamma_m\backslash\mathsf G_m}\sum_{0\ne y\in\mathcal O^{1\times m}}f(y\gamma)\,d\mu_m=\int f\,dx\).

\textbf{(U) {[}L, from memory; Bourbaki Int. VII §2 Th. 3, Folland AHA Thm 2.49{]}} A \(\mathsf G\)-invariant Radon measure on a homogeneous space \(\mathsf H\backslash\mathsf G\) (\(\mathsf H\) closed) is unique up to a scalar.

\textbf{(Tr) {[}P, = R17 (R2) transposed{]}} For \(m\ge2\), \(\mathsf G_m\) is transitive on \(\Omega_m\) (at each place move \(x_v\) to \(y_v\) by \(h_v\in GL_m(\mathbb K_v)\) and correct \(|\det_{\mathbb R}h_v|\) to 1 inside the stabiliser of \(y_v\), which contains conjugates of \(\mathrm{diag}(1,c,1,..)\)); the orbit map \(\mathsf H_y\backslash\mathsf G_m\to\Omega_m\) is a homeomorphism {[}L, standard{]}. Every \(0\ne y\in\Delta^{1\times m}\) lies in \(\Omega_m\) (a nonzero coordinate is invertible in \(\Delta\), so nonzero at every place) {[}P{]}.

\begin{center}\rule{0.5\linewidth}{0.5pt}\end{center}

\hypertarget{part-a.-rank-one-formula-for-commensurable-lattices}{%
\subsubsection*{PART A. Rank-one formula for commensurable lattices}\label{part-a.-rank-one-formula-for-commensurable-lattices}}

\textbf{Theorem A {[}P from (L1)\(_m\), (U){]}.} Let \(m\ge2\), \(L:=\mathcal O^{1\times m}\), \(N\subset\Delta_{\mathbb R}^{1\times m}\) a lattice commensurable with \(L\) (any subgroup, no module structure required), \(\Gamma'\subseteq\mathsf G_m\) a subgroup commensurable with \(GL_m(\mathcal O)\) with \(N\Gamma'=N\). Then \(\Gamma'\) is a lattice in \(\mathsf G_m\) and, with \(\mu'\) the invariant probability measure on \(\Gamma'\backslash\mathsf G_m\), for every Borel \(f:\Delta_{\mathbb R}^{1\times m}\to[0,\infty]\)
\[\int_{\Gamma'\backslash\mathsf G_m}\sum_{0\ne y\in N}f(y\gamma)\,d\mu'(\gamma)=\frac1{\mathrm{covol}(N)}\int f(x)\,dx .\]

\emph{Proof.} \textbf{A.1 (congruence subgroup).} Choose an integer \(a\ge1\) with \(aL\subseteq N\subseteq a^{-1}L\). Put \(\Gamma'':=\{\delta\in\Gamma'\cap GL_m(\mathcal O):\delta\equiv1\bmod a^2M_m(\mathcal O)\}\): the kernel of \(\Gamma'\cap GL_m(\mathcal O)\to GL_m(\mathcal O/a^2\mathcal O)\) (finite), so of finite index in \(\Gamma'\cap GL_m(\mathcal O)\), hence in \(\Gamma'\) and in \(\Gamma_m\). For \(c\in a^{-1}L\), \(\delta\in\Gamma''\): \(c\delta-c\in a^{-1}L\cdot a^2M_m(\mathcal O)=aL\). So every coset \(c+aL\subseteq a^{-1}L\) is \(\Gamma''\)-stable. \(\Gamma''\) is a lattice in \(\mathsf G_m\) (finite index in \(\Gamma_m\)); \(\Gamma'\) is discrete (finite union of cosets of the discrete closed \(\Gamma''\)) and contains \(\Gamma''\) with finite index, so it is a lattice.

\textbf{A.2 (finite covers).} If \(\Gamma''\le\Gamma'\) has index \(\ell\) and \(\mathcal F'\) is a Borel fundamental domain for \(\Gamma'\), then \(\bigsqcup_{\delta\in\Gamma''\backslash\Gamma'}\delta\mathcal F'\) is one for \(\Gamma''\); so for a left \(\Gamma'\)-invariant Borel \(\Phi\ge0\) the mean over \(\Gamma''\backslash\mathsf G_m\) equals the mean over \(\Gamma'\backslash\mathsf G_m\) (both for the probability normalisation). Same for \(\Gamma''\le\Gamma_m\).

\textbf{A.3 (coset measures).} For \(c\in a^{-1}L\) and Borel \(B\): \(\nu_c(B):=\int_{\Gamma''\backslash\mathsf G_m}\#\big(((c+aL)\gamma\cap B)\setminus\{0\}\big)\,d\mu''(\gamma)\). The integrand is Borel and left \(\Gamma''\)-invariant (A.1), \(\nu_c\) is a Borel measure (monotone convergence) and \(\nu_c(Bg)=\nu_c(B)\) for \(g\in\mathsf G_m\) (right invariance of \(\mu''\)). Summing over the \([a^{-1}L:aL]=a^{2Dm}\) cosets and using A.2, (L1)\(_m\):
\[\sum_{c\in a^{-1}L/aL}\nu_c(B)=\int_{\Gamma''\backslash\mathsf G_m}\#(a^{-1}L\gamma\cap B\setminus0)\,d\mu''=\int_{\Gamma_m\backslash\mathsf G_m}\#(L\gamma\cap aB\setminus0)\,d\mu_m=\mathrm{Leb}(aB)=a^{Dm}\mathrm{Leb}(B).\tag{A.3}\]
Hence \(\nu_c\le a^{Dm}\mathrm{Leb}\): \(\nu_c\) is Radon and \(\nu_c(\Delta_{\mathbb R}^{1\times m}\setminus\Omega_m)=0\) (a Lebesgue null set; alternatively: every nonzero point of \((c+aL)\gamma\) lies in \(\Omega_m\) by (Tr), as \(c+aL\subset\Delta^{1\times m}\) --- the check asked for in the task is positive). By (Tr), (U) applied to \(\nu_c|_{\Omega_m}\) and \(\mathrm{Leb}|_{\Omega_m}\) (invariant since \(\|g\|=1\)): \(\nu_c=C_c\,\mathrm{Leb}\), \(0\le C_c\le a^{Dm}\).

\textbf{A.4 (the constant, averaging trick; no domination needed).} Let \(0\le f\in C_c\), \(\int f>0\), \(f_\varepsilon(x)=f(\varepsilon x)\). Then \(\varepsilon^{Dm}\int f_\varepsilon\,d\nu_c=C_c\int f\) for all \(\varepsilon>0\). For every \(\gamma\), \((c+aL)\gamma\) is a translate of the lattice \(aL\gamma\) of covolume \(a^{Dm}\), so \(\varepsilon^{Dm}\sum_{y\in(c+aL)\gamma\setminus0}f(\varepsilon y)\to a^{-Dm}\int f\) as \(\varepsilon\to0\) (Riemann sums of a uniformly continuous compactly supported function over a translated lattice; the omitted term \(\varepsilon^{Dm}f(0)\to0\)). Fatou: \(a^{-Dm}\int f\le\liminf_\varepsilon\varepsilon^{Dm}\int f_\varepsilon d\nu_c=C_c\int f\). So \(C_c\ge a^{-Dm}\) for EVERY coset \(c\). By (A.3), \(\sum_cC_c=a^{Dm}\) over \(a^{2Dm}\) cosets, each term \(\ge a^{-Dm}\): hence \(C_c=a^{-Dm}=\mathrm{covol}(aL)^{-1}\) for all \(c\in a^{-1}L\). (Inputs used: (L1)\(_m\) for \(\mathcal O^{1\times m}\) and Borel sets, (U), (Tr), Fatou. Gargava\textquotesingle s Lemma 18 is used only inside (L1).)

\textbf{A.5 (conclusion).} \(N\setminus0=\bigsqcup_{c\in N/aL}(c+aL)\setminus0\); by A.2 and monotone convergence the left side of the theorem is \(\sum_{c\in N/aL}\int f\,d\nu_c=[N:aL]\,a^{-Dm}\int f=\mathrm{covol}(N)^{-1}\int f\). \(\blacksquare\)

\textbf{Lemma A.6 (stabilisers are commensurable) {[}P{]}.} If \(L_0\subset\Delta^{1\times m}\) is a lattice commensurable with \(\mathcal O^{1\times m}\) then \(\mathrm{Stab}_{GL_m(\Delta_{\mathbb R})}(L_0)\subseteq\mathsf G_m\) and it is commensurable with \(GL_m(\mathcal O)\).
\emph{Proof.} A stabiliser of a lattice has \(|\det_{\mathbb R}|=1\). With \(aL\subseteq L_0\subseteq a^{-1}L\): the congruence subgroup \(\Gamma_m(a^2)\) fixes all cosets of \(aL\) in \(a^{-1}L\), so it stabilises \(L_0\). Conversely the finite-index subgroup of \(\mathrm{Stab}(L_0)\) acting trivially on the finite group \(a^{-1}L_0/aL_0\) fixes every coset of \(aL_0\) in \(a^{-1}L_0\), hence stabilises \(L\) (\(aL_0\subseteq L\subseteq a^{-1}L_0\)), hence lies in \(GL_m(\mathcal O)\) (\(L\gamma=L\Rightarrow\gamma,\gamma^{-1}\in M_m(\mathcal O)\)). \(\square\)

\begin{center}\rule{0.5\linewidth}{0.5pt}\end{center}

\hypertarget{part-b.-j-frames-for-arbitrary-module-lattices}{%
\subsubsection*{\texorpdfstring{PART B. \(j\)-frames for arbitrary module lattices}{PART B. j-frames for arbitrary module lattices}}\label{part-b.-j-frames-for-arbitrary-module-lattices}}

\textbf{Notation.} \(1\le j<k\). \(L_0\subset\Delta^{1\times k}\) a full left \(\mathcal O\)-lattice; \(M\subset\Delta^{j\times1}\) a full right \(\mathcal O\)-lattice; \(ML_0\subset\Delta^{j\times k}\) the subgroup generated by products (column)(row); it is a lattice in \(\Delta_{\mathbb R}^{j\times k}\) (between \(a^2\mathcal O^{j\times k}\) and \(a^{-2}\mathcal O^{j\times k}\)). \(\Gamma:=\mathrm{Stab}_{\mathsf G_k}(L_0)\) (commensurable with \(GL_k(\mathcal O)\) by A.6, a lattice in \(\mathsf G_k\) by Theorem A), \(\mu\) the invariant probability measure on \(\Gamma\backslash\mathsf G_k\). \((ML_0)\gamma=ML_0\) for \(\gamma\in\Gamma\).

\textbf{Theorem B {[}P from (Her), Theorem A in ranks \(k,\dots,k-j+1\), i.e. from (L1)\(_m\), \(k-j+1\le m\le k\){]}.} For every Borel \(F:\Delta_{\mathbb R}^{j\times k}\to[0,\infty]\)
\[\int_{\Gamma\backslash\mathsf G_k}\ \sum_{Y\in ML_0,\ \mathrm{rk}_\Delta Y=j}F(Y\gamma)\,d\mu(\gamma)=\frac1{\mathrm{covol}(ML_0)}\int_{\Delta_{\mathbb R}^{j\times k}}F(y)\,dy\qquad\text{in }[0,\infty].\tag{B$_{j,k}(M,L_0)$}\]

\emph{Proof.} Induction on \(j\), simultaneously for all \(k>j\) and all pairs \((M,L_0)\).

\textbf{B.0 Base \(j=1\).} \(M=I_1\) is a full right \(\mathcal O\)-lattice in \(\Delta\), \(N_1:=I_1L_0\) is a \(\Gamma\)-stable lattice commensurable with \(\mathcal O^{1\times k}\), rank 1 means \(\ne0\): Theorem A for \((N_1,\Gamma)\), \(m=k\ge2\).

\textbf{B.1 Splitting \(M\) (uses (Her), right modules).} Let \(j\ge2\) (so \(k\ge3\)). \(\pi_1:\Delta^{j\times1}\to\Delta\) the first coordinate (right \(\Delta\)-linear), \(I_1:=\pi_1(M)\), a full right \(\mathcal O\)-lattice in \(\Delta\), projective by (Her). So \(\pi_1|_M\) has a right \(\mathcal O\)-linear section, which extends \(\Delta\)-linearly: \(x\mapsto sx\) with \(s=(1,s'')^T\in\Delta^{j\times1}\), \(sI_1\subseteq M\), and \(M=sI_1\oplus(0\times M_1)\), \(0\times M_1:=M\cap\ker\pi_1\), \(M_1\) a full right \(\mathcal O\)-lattice in \(\Delta^{(j-1)\times1}\). Let \(A=\begin{pmatrix}1&0\\-s''&1\end{pmatrix}\in GL_j(\Delta)\): \(As=e_1\), \(AM=e_1I_1\oplus(0\times M_1)\). \(Y\mapsto AY\) commutes with the right action of \(\mathsf G_k\), preserves \(\mathrm{rk}_\Delta\) and Lebesgue measure (unipotent), and maps \(ML_0\) onto \((AM)L_0\). So (B\(_{j,k}(M,L_0)\)) for \(F\) is (B\(_{j,k}(AM,L_0)\)) for \(F\circ A^{-1}\); \textbf{WLOG \(M=e_1I_1\oplus(0\times M_1)\).} Then
\[ML_0=\Big\{\tbinom{y}{Y_1}:\ y\in N_1:=I_1L_0,\ Y_1\in M_1L_0\Big\},\qquad \mathrm{covol}(ML_0)=\mathrm{covol}(N_1)\,\mathrm{covol}(M_1L_0)\]
(the generators \((e_1x+\binom0{m_1})l\) and the elements \(e_1x\), \(\binom0{m_1}\) of \(M\) generate the same group). \(N_1\) is a \(\Gamma\)-stable lattice in \(\Delta_{\mathbb R}^{1\times k}\) (not a left \(\mathcal O\)-module in general), \(M_1L_0\) is \(\Gamma\)-stable.

\textbf{B.2 Orbits and unfolding.} \(N_1\setminus0=\bigsqcup_{i\in I}y_i\Gamma\) (countable). \(\mathsf H_i:=\mathrm{Stab}_{\mathsf G_k}(y_i)\), \(\Gamma_i:=\Gamma\cap\mathsf H_i\). As in {[}R15, §3 Step 1{]}:
\[S_jF(\gamma):=\sum_{Y\in ML_0,\mathrm{rk}Y=j}F(Y\gamma)=\sum_i\sum_{\delta\in\Gamma_i\backslash\Gamma}\Phi_i(\delta\gamma),\quad\Phi_i(g):=\sum_{Y_1\in M_1L_0,\ \mathrm{rk}\binom{y_i}{Y_1}=j}F\tbinom{y_ig}{Y_1g},\]
\(\Phi_i\) Borel, left \(\Gamma_i\)-invariant. By B.3 below \(\mathsf H_i\) is unimodular and \(\Gamma_i\) is a lattice in it; let \(\eta_i\) be the Haar measure of \(\mathsf H_i\) with \(\eta_i(\Gamma_i\backslash\mathsf H_i)=1\) and \(\nu_i\) the invariant measure on \(\mathsf H_i\backslash\mathsf G_k\) of {[}R15, Lemma 2.6{]} ({[}L{]} Weil\textquotesingle s quotient formula). Then
\[\int_{\Gamma\backslash\mathsf G_k}S_jF\,d\mu=\sum_i\int_{\mathsf H_i\backslash\mathsf G_k}\Big[\int_{\Gamma_i\backslash\mathsf H_i}\Phi_i(hg)\,d\eta_i(h)\Big]d\nu_i(g).\tag{B.2}\]

\textbf{B.3 Normal form at the orbit \(i\) (uses (Her), left modules).} \(y_i\in\Delta^{1\times k}\setminus0\); choose \(g_i\in GL_k(\Delta)\) with \(e_1g_i=y_i\). Since \(\mathsf G_k\) is NORMAL in \(GL_k(\Delta_{\mathbb R})\), \(g_i\mathsf H_ig_i^{-1}=\mathrm{Stab}_{\mathsf G_k}(e_1)=\mathsf H:=\{u(n,h)=\begin{pmatrix}1&0\\n&h\end{pmatrix}:n\in\Delta_{\mathbb R}^{(k-1)\times1},h\in\mathsf G_{k-1}\}\) (no condition on \(\|g_i\|\) is needed; \(\|u(n,h)\|=\|h\|\)). \(L_i:=L_0g_i^{-1}\) is a full left \(\mathcal O\)-lattice and \(g_i\Gamma_ig_i^{-1}=\mathrm{Stab}_{\mathsf H}(L_i)\). Let \(p:\Delta^{1\times k}\to\Delta^{1\times(k-1)}\) forget the first coordinate; \(L_i'':=p(L_i)\) is a full left \(\mathcal O\)-lattice, projective, so \(p|_{L_i}\) has a section \(x''\mapsto(x''t\,|\,x'')\), \(t\in\Delta^{(k-1)\times1}\), and \(L_i=(Je_1\oplus(0|L_i''))\,u(t,1)\) with \(Je_1:=L_i\cap\Delta e_1\), \(J\) a full left \(\mathcal O\)-lattice in \(\Delta\). Replacing \(g_i\) by \(u(t,1)g_i\) (still \(e_1g_i=y_i\)): \textbf{\(L_i=Je_1\oplus(0\,|\,L_i'')\).} Then, from \((a|x'')u(n,h)=(a+x''n\,|\,x''h)\) and \(u(n,h)^{-1}=u(-h^{-1}n,h^{-1})\):
\[\mathrm{Stab}_{\mathsf H}(L_i)=\{u(n,h):\ h\in\Gamma_i'':=\mathrm{Stab}_{\mathsf G_{k-1}}(L_i''),\ n\in T_i\},\qquad T_i:=\{n:\ L_i''n\subseteq J\}.\]
\(T_i\) is a lattice in \(\Delta_{\mathbb R}^{(k-1)\times1}\) (if \(c\mathcal O^{1\times(k-1)}\subseteq L_i''\subseteq c^{-1}\mathcal O^{1\times(k-1)}\) and \(c\mathcal O\subseteq J\subseteq c^{-1}\mathcal O\), \(c\in\mathbb Z_{>0}\), then \(c^2\mathcal O^{(k-1)\times1}\subseteq T_i\subseteq c^{-2}\mathcal O^{(k-1)\times1}\)) and \(hT_i=T_i\) for \(h\in\Gamma_i''\). \(\Gamma_i''\) is a lattice in \(\mathsf G_{k-1}\) (A.6 + Theorem A, rank \(k-1\ge2\)); let \(\mu_i''\) be the Haar measure of \(\mathsf G_{k-1}\) with \(\mu_i''(\Gamma_i''\backslash\mathsf G_{k-1})=1\). As in {[}R15, Lemma 2.2{]}: \(\mathsf H\cong\mathsf G_{k-1}\ltimes\Delta_{\mathbb R}^{(k-1)\times1}\) is unimodular with Haar measure \(dn\,dh\), and \(\{u(n,h):h\in\mathcal F_i'',n\in\mathcal P_i\}\) (\(\mathcal F_i''\) a fundamental domain of \(\Gamma_i''\), \(\mathcal P_i\) a fundamental parallelotope of \(T_i\)) is a fundamental domain for \(\mathrm{Stab}_{\mathsf H}(L_i)\backslash\mathsf H\) (from \(u(m,\sigma)u(n,h)=u(m+\sigma n,\sigma h)\)). So \(\mathrm{Stab}_{\mathsf H}(L_i)\) is a lattice in \(\mathsf H\), the normalised Haar measure is \(\mathrm{covol}(T_i)^{-1}dn\,d\mu_i''(h)\), and conjugation by \(g_i\) transports all this to \((\mathsf H_i,\Gamma_i,\eta_i)\).
Moreover \(M_1L_i=\{(b\,|\,Y''):b\in M_1J,\ Y''\in M_1L_i''\}\), \(M_1J\subset\Delta^{(j-1)\times1}\) a lattice, and \(\mathrm{covol}(M_1L_i)=\mathrm{covol}(M_1J)\,\mathrm{covol}(M_1L_i'')\).

\textbf{B.4 Torus lemma.} Let \(Y''\in M_1L_i''\), \(\mathrm{rk}_\Delta Y''=j-1\). Then \(Y''T_i\subseteq M_1J\) (write \(Y''=\sum_\ell m_\ell x_\ell''\), \(m_\ell\in M_1\), \(x_\ell''\in L_i''\); \(Y''n=\sum m_\ell(x_\ell''n)\)), and \(n\mapsto Y''n\) induces a continuous surjective homomorphism of compact groups \(\Delta_{\mathbb R}^{(k-1)\times1}/T_i\to\Delta_{\mathbb R}^{(j-1)\times1}/M_1J\) (\(Y''\) has a right inverse over \(\Delta\)). It maps normalised Haar measure to normalised Haar measure, so for Borel \(\psi\ge0\):
\[\frac1{\mathrm{covol}(T_i)}\int_{\mathcal P_i}\sum_{b\in M_1J}\psi(b+Y''n)\,dn=\frac1{\mathrm{covol}(M_1J)}\int_{\Delta_{\mathbb R}^{(j-1)\times1}}\psi(t)\,dt.\]

\textbf{B.5 Fibre mean value.} Claim: for Borel \(\Psi\ge0\) on \(\Delta_{\mathbb R}^{(j-1)\times k}\),
\[\int_{\Gamma_i\backslash\mathsf H_i}\sum_{Y_1\in M_1L_0,\ \mathrm{rk}\binom{y_i}{Y_1}=j}\Psi(Y_1h)\,d\eta_i(h)=\frac1{\mathrm{covol}(M_1L_0)}\int\Psi(Z)\,dZ\qquad\text{(independent of }i).\]
Substitute \(h=g_i^{-1}h_1g_i\), \(Y_1=Y'g_i\), \(Y'=(b|Y'')\in M_1L_i\): \(\mathrm{rk}\binom{y_i}{Y_1}=\mathrm{rk}\binom{e_1}{Y'}=1+\mathrm{rk}Y''\), \(Y_1h=Y'h_1g_i\), \(Y'u(n,h)=(b+Y''n\,|\,Y''h)\). With \(\Psi^{g_i}(Z):=\Psi(Zg_i)\), B.3, Tonelli and B.4:
\[\text{LHS}=\int_{\mathcal F_i''}\sum_{Y''\in M_1L_i'',\ \mathrm{rk}=j-1}\Big[\frac1{\mathrm{covol}T_i}\int_{\mathcal P_i}\sum_{b\in M_1J}\Psi^{g_i}(b+Y''n\,|\,Y''h)dn\Big]d\mu_i''(h)=\frac1{\mathrm{covol}(M_1J)}\int_{\Gamma_i''\backslash\mathsf G_{k-1}}\sum_{Y''}\Psi_1(Y''h)\,d\mu_i'',\]
\(\Psi_1(Z''):=\int\Psi^{g_i}(t|Z'')dt\). By the induction hypothesis (B\(_{j-1,k-1}(M_1,L_i'')\)) (\(1\le j-1<k-1\)) this is \(\mathrm{covol}(M_1J)^{-1}\mathrm{covol}(M_1L_i'')^{-1}\int\Psi_1=\mathrm{covol}(M_1L_i)^{-1}\int\Psi(Zg_i)dZ\). Jacobians: \(\int\Psi(Zg_i)dZ=\|g_i\|^{-(j-1)}\int\Psi\) and \(\mathrm{covol}(M_1L_i)=\mathrm{covol}((M_1L_0)g_i^{-1})=\|g_i\|^{-(j-1)}\mathrm{covol}(M_1L_0)\); they cancel. \(\square\)

\textbf{B.6 Conclusion.} In (B.2), \(\Phi_i(hg)=\sum_{Y_1}\Psi_g(Y_1h)\), \(\Psi_g(Z)=F\binom{y_ig}{Zg}\); by B.5 and \(\|g\|=1\) the inner integral is \(\mathrm{covol}(M_1L_0)^{-1}f(y_ig)\), \(f(x):=\int F\binom xZdZ\). The same unfolding (B.2) for \(j=1\), the lattice \(N_1\) and the function \(f\) gives \(\int_{\Gamma\backslash\mathsf G_k}\sum_{0\ne y\in N_1}f(y\gamma)d\mu=\sum_i\int_{\mathsf H_i\backslash\mathsf G_k}f(y_ig)d\nu_i(g)\) with the same \(\nu_i\). Theorem A for \((N_1,\Gamma)\) in rank \(k\): this is \(\mathrm{covol}(N_1)^{-1}\int f\). Hence \(\int S_jF\,d\mu=\mathrm{covol}(M_1L_0)^{-1}\mathrm{covol}(N_1)^{-1}\int F=\mathrm{covol}(ML_0)^{-1}\int F\). \(\blacksquare\)

\textbf{Rank-one inputs used.} Theorem A in rank \(k\) for \((I_1L_0,\mathrm{Stab}(L_0))\); in rank \(k-1\) for \((I_2L_i'',\mathrm{Stab}(L_i''))\) for every orbit \(i\) (and the lattice property of \(\mathrm{Stab}(L_i'')\)); ...; down to rank \(k-j+1\ge2\). Through Theorem A these reduce to (L1)\(_m\) for \((\mathcal O^{1\times m},GL_m(\mathcal O))\), \(m=k-j+1,\dots,k\) only. For \(M=\mathcal O^{j\times1}\), \(L_0=\mathcal O^{1\times k}\) this is Theorem R of R15 \textbf{without (h1\(^\ell\))} (for hereditary \(\mathcal O\)).

\begin{center}\rule{0.5\linewidth}{0.5pt}\end{center}

\hypertarget{part-c.-theorem-e-without-h1}{%
\subsubsection*{PART C. Theorem E without (h1)}\label{part-c.-theorem-e-without-h1}}

\textbf{C.1 Where (h1), (h1\(^\ell\)) were used {[}read: R14 §0-1, R15, R18{]}.} (i) R14 0.1: coordinates \(\Lambda_W=\mathcal O^r\) (Lemma 1.1). (ii) R14 Lemma 1.2: \(\mathcal O^r\cap U=B_U\mathcal O^j\) free. (iii) R15 Lemma 2.1 (h1\(^\ell\)): \(y=ae_1\delta\), i.e. saturations of rows are free summands with free complement. Nowhere else in the parts read.

\textbf{C.2 Coordinates without freeness {[}P{]}.} Choose any \(\Delta\)-basis of \(W\): \(W=\Delta^{r\times1}\), \(\Lambda_W\subset\Delta^{r\times1}\) a full right \(\mathcal O\)-lattice, \(G\subset GL_r(\Delta)\) acting from the left with \(g\Lambda_W=\Lambda_W\). R14 Lemma 0.1 (\(M_v\in GL_r(\mathbb K)\), \(u_g=MgM^{-1}\) unitary) uses only \(W_{\mathbb R}=\bigoplus_v\mathbb K^r\): unchanged. For \(L=\mathcal O^{1\times k}\gamma\): \(\Lambda_W\otimes_{\mathcal O}L\cong\Lambda_W\mathcal O^{1\times k}\gamma=\Lambda_W^{\,k}\gamma\) (\(\Lambda_W^{\,k}\):= matrices in \(\Delta^{r\times k}\) with all columns in \(\Lambda_W\); \(L\) is free so the tensor product is \(\Lambda_W^k\)), \(\Lambda(L):=M\Lambda_W^{\,k}\gamma\subset E\), \(\mathrm{covol}(M\Lambda_W^{\,k})=c_W^k\), \(u_g\Lambda(L)=\Lambda(L)\). R14 Lemmas 0.2-0.4 are statements about \(E\) and \(\mathcal Q\) only. Lemma 0.5 holds with \(X\in\Lambda_W^{\,k}\) (same proof: \(gX=X\)).

\textbf{C.3 Lemma 1.2\textquotesingle{} (MVF without (h1)) {[}P from Theorem B{]}.} Let \(\mathcal O\) be maximal, \(k>r\), \(\lambda=(\kappa_n/c_W^k)^{1/n}\), \(\gamma\) distributed by \(\mu_k\) on \(\Gamma_k\backslash\mathsf G_k\), \(\Lambda=\lambda M\Lambda_W^{\,k}\gamma\). Then \(\mathrm{covol}(\Lambda)=\kappa_n\) and for Borel \(\varphi\ge0\), (MVF) of R14 holds verbatim with the same \(\mu_U=\kappa_n^{-f_U}(c_W^{f_U}/c_U)^k\mathrm{vol}_{E_U}\), including the "fixed \(U\)" statement.
\emph{Proof.} For \(0\ne X\in\Lambda_W^{\,k}\) let \(U=X\Delta^{k\times1}\), \(j=\dim_\Delta U\). \(\{X\in\Lambda_W^{\,k}:\mathrm{col}(X)=U\}=\{X\in(\Lambda_W\cap U)^k:\mathrm{rk}_\Delta X=j\}\). Fix \(B_U\in\Delta^{r\times j}\) whose columns are a \(\Delta\)-basis of \(U\) and put \(M_U:=\{m\in\Delta^{j\times1}:B_Um\in\Lambda_W\}\), a full right \(\mathcal O\)-lattice (\(\Lambda_W\cap U\) is a full \(\mathbb Z\)-lattice in the \(\mathbb Q\)-subspace \(U\) and a right \(\mathcal O\)-module). Then \((\Lambda_W\cap U)^k=B_U(M_U\mathcal O^{1\times k})\), \(\mathrm{rk}(B_UY)=\mathrm{rk}\,Y\). Theorem B for \((M_U,L_0=\mathcal O^{1\times k})\), \(\Gamma=GL_k(\mathcal O)\), \(F(y)=\varphi(\lambda MB_Uy)\) (\(j\le r<k\)): the \(U\)-part has mean \(\mathrm{covol}(M_U\mathcal O^{1\times k})^{-1}\int F\,dy\). The linear bijection \(y\mapsto\lambda MB_Uy:\Delta_{\mathbb R}^{j\times k}\to E_U\) maps \(M_U\mathcal O^{1\times k}\) onto \(\lambda(M(\Lambda_W\cap U))^k\), of Euclidean covolume \(\lambda^{fn}c_U^k\); for a linear isomorphism \(T\) and a lattice \(N\), \(\mathrm{covol}(N)^{-1}\int\varphi\circ T=\mathrm{covol}(TN)^{-1}\int\varphi\). So the mean is \((\lambda^{fn}c_U^k)^{-1}\int_{E_U}\varphi=\int\varphi\,d\mu_U\); \(U=W\) gives \(\kappa_n^{-1}\int_E\varphi\). \(\square\)
Lemma 1.4 (formula for \(T_U(V)\), bound by \(R_U^{k/2}\)) and Lemma 1.5 (counting, Minkowski) involve only \(c_U,c_W\) and \(\Lambda_W\) as a \(\mathbb Z\)-lattice: unchanged.

\textbf{C.4 Rest of R14 {[}read in this session, all sections, for this question only{]}.}

\begin{itemize}
\tightlist
\item
  §2: Prop. 2.1 {[}V from R6{]} is about an arbitrary lattice and \(\mathcal Q\ni\mathrm{Id}\); Lemma 2.2 uses \(u_g\Lambda=\Lambda\), \(A_t\in\mathcal Q\) and Lemma 0.3(a) only ("non-free" there refers to the \(G\)-action on lattice points, not to \(\mathcal O\)-modules); Prop. 2.3 uses \(0\le N_t\le D_k\).
\item
  §3: Prop. 3.1, 3.2, Lemma 3.3, 3.4, Prop. 3.5, Cor. 3.6 concern single points, the spaces \(E_U\), the function \(q\) (Lemma 0.4: \(Q_v\) = orthonormal basis of \(M_vU_v\), no lattice), and \(T_U\), \(R_U\), \(\Sigma_k\) (Lemmas 1.4, 1.5: only \(c_U\), \(c_W\), Minkowski for the \(\mathbb Z\)-lattice \(\Lambda_W\)).
\item
  §4: Prop. 4.1 uses (MVF), its fixed-\(U\) version, Lemma 0.5 (\(\Lambda^{nf}\subseteq\bigcup_{g\ne1}\{\mathrm{col}(x)\subseteq W^g\}\)) and Markov on ONE probability space: here \((\Gamma_k\backslash\mathsf G_k,\mu_k)\) for every \(U\), because \(L_0=\mathcal O^{1\times k}\) and \(\mathrm{Stab}_{\mathsf G_k}(L_0)=GL_k(\mathcal O)\) in every application of Theorem B. Read "\(\lambda M\mathcal O^{r\times k}\gamma\)" as \(\lambda M\Lambda_W^{\,k}\gamma\).
\item
  §5: the final step replaces \(\gamma\) by \(\gamma N\), \(N\in GL_k(\Delta_{\mathbb R})\) positive Hermitian: \(\sqrt{A_T}\Lambda=\lambda M\Lambda_W^{\,k}\gamma N\) stays in the family.
\item
  §6 (Leech) is not used for Theorem E\textquotesingle.
  So freeness of \(\Lambda_W\) and (h1) enter R14 only through §0.1 (coordinates), Lemma 1.1 and the factorisation \(X=B_UY\), \(B_U\in\mathcal O^{r\times j}\), \(Y\in\mathcal O^{j\times k}\) in Lemma 1.2; (h1\(^\ell\)) only through R15 Lemma 2.1. Both are replaced above.
\end{itemize}

\textbf{Theorem E\textquotesingle{} (no class number hypothesis).} DATA (T) as in R18 (\(G\), \(W\), \(\Delta=\mathrm{End}_G(W)^{op}\), \(\Delta_{\mathbb R}\cong\mathbb K^s\), \(\Lambda_W\) a \(G\)-invariant lattice, \(\mathcal O=\{c:\Lambda_Wc\subseteq\Lambda_W\}\)). HYPOTHESES: \textbf{\(\mathcal O\) is a maximal order} (only (Her) is used); \textbf{(BS)} for all proper nonzero \(\Delta\)-subspaces \(U\). INPUTS: (L1)\(_m\), \(m\ge2\), for \((\mathcal O^{1\times m},GL_m(\mathcal O),\mathsf G_m)\) {[}= Gargava for \((\Delta^{op},\mathcal O^{op})\){]}; (L2); (L3); (U); (Her). CONCLUSION: there is \(k_0\) (non-effective) such that for \(k\ge k_0\), \(n=kN_W\), there is \(\gamma'\in GL_k(\Delta_{\mathbb R})\) such that the \(G\)-invariant lattice \(\Lambda_W\otimes_{\mathcal O}\mathcal O^{1\times k}\gamma'=M\Lambda_W^{\,k}\gamma'\) has an empty open ball of relative volume \((1-o(1))\,|G|\,n^2/(4e\dim_{\mathbb Q}A)\). Hence \(V_n\ge(1-o(1))|G|n^2/(4e\dim_{\mathbb Q}A)\) for \(N_W\mid n\).
\emph{Proof.} R14 §0-5 with C.2, Lemma 1.2\textquotesingle{} in place of Lemma 1.1/1.2, and Theorem B in place of Theorem R. \(\square\) (Status: {[}P{]} modulo C.4.)
Remark: if \(\mathcal O(\Lambda_W)\) is not maximal, \(\Lambda_W\mathcal O'\) (\(\mathcal O'\supseteq\mathcal O\) maximal) is a \(G\)-invariant lattice with maximal right order; (BS) must then be checked for that lattice.

\textbf{Corollary (Weil lattices \(W_p\) of \(SL_2(p)\)) {[}conditional{]}.} Let \(p\equiv5\bmod8\), \((p+1)/2\) prime, and take the data recorded in earlier rounds {[}L, from the task statement; NOT re-derived here{]}: \(|G|=p(p^2-1)\), \(N_W=2(p-1)\), \(\Delta_{\mathbb R}=\mathbb H\times\mathbb H\), \(r=(p-1)/4\), \(\dim_{\mathbb Q}A=r^2sd=(p-1)^2/2\). IF \(\mathcal O(W_p)\) is maximal and \(W_p\) satisfies (BS) for ALL proper \(\Delta\)-subspaces, then \(V_n\ge(1-o(1))\frac{p(p+1)}{2e(p-1)}n^2\) for \(2(p-1)\mid n\), \(n\to\infty\), with no class number hypothesis (\(p(p^2-1)/(4e(p-1)^2/2)=p(p+1)/(2e(p-1))\); \(p=13\): \(91/(12e)\)). (BS) is open for \(p\ge61\) and false for \(p=37\) {[}as stated by the orchestrator; not checked here{]}; maximality of \(\mathcal O(W_p)\) must be verified {[}not checked here{]}.

\begin{center}\rule{0.5\linewidth}{0.5pt}\end{center}

\hypertarget{d.-status-gaps-what-a-referee-should-check}{%
\subsubsection*{D. Status, gaps, what a referee should check}\label{d.-status-gaps-what-a-referee-should-check}}

\begin{enumerate}
\def\labelenumi{\arabic{enumi}.}
\tightlist
\item
  \textbf{Part A} complete {[}P{]}. Check: A.3 (the sum over ALL cosets of \(aL\) in \(a^{-1}L\) is the Siegel transform of \(a^{-1}L\), which is \(\Gamma_m\)-invariant); A.4 (Fatou lower bound for every coset + exact total \(\Rightarrow\) equality); that (U) is applied to Radon measures on \(\Omega_m\cong\mathsf H_y\backslash\mathsf G_m\).
\item
  \textbf{Part B} complete {[}P{]}. Check: B.1 (the section of \(\pi_1\) is \(x\mapsto sx\) with \(\pi_1(s)=1\); \(ML_0\) is a product lattice); B.3 (description of \(\mathrm{Stab}_{\mathsf H}(L_i)\), \(T_i\) a lattice, normality of \(\mathsf G_k\) makes any \(g_i\in GL_k(\Delta)\) admissible); B.5 (Jacobians \(\|g_i\|^{\pm(j-1)}\) cancel; normalisation \(\mathrm{covol}(T_i)^{-1}dn\,d\mu_i''\)); B.6 (same \(\nu_i\) for \(j\) and for 1).
\item
  \textbf{(Her)} is cited from memory (Reiner). It is the only place where maximality is used; the statements hold for any hereditary order.
\item
  \textbf{(L1)\(_m\)} is used as in R18 (family \(\mathsf G_m=\{\|\cdot\|=1\}\), \(\Gamma_m=GL_m(\mathcal O)\)); its status is that of R17 (Weil\textquotesingle s reduction theory unread). If some \(\mathbb K_v=\mathbb R\), R18 problem 6 applies unchanged.
\item
  \textbf{Part C}: Lemma 1.2\textquotesingle{} is complete {[}P{]}. Theorem E\textquotesingle{} = R14 §0-§5 with the substitutions of C.2-C.4; the analytic part of R14 (§2-§5) was read here only to locate uses of freeness (none), not re-verified (R18 did not re-verify it either; R14\textquotesingle s own referee reports were not read here).
\item
  Things a referee should try to break: (a) A.4: is \([a^{-1}L:aL]=a^{2Dm}\) and \(\mathrm{covol}(aL)=a^{Dm}\) with \(D=\dim_{\mathbb Q}\Delta\) (yes: \(L\cong\mathbb Z^{Dm}\)); (b) B.3 when \(J\ne\mathcal O\) or \(L_i''\) non-free: nothing in B.3-B.5 uses freeness; (c) that the induction hypothesis is applied to the group \(\mathrm{Stab}_{\mathsf G_{k-1}}(L_i'')\), which varies with the orbit \(i\) (allowed: the induction is over all pairs); (d) C.3: the same probability space for all \(U\).
\item
  Corollary for \(W_p\): purely conditional; the data of \(W_p\), maximality of \(\mathcal O(W_p)\) and (BS) are not verified here.
\item
  By-products {[}P{]}: (i) affine/congruence Siegel formula: for every coset, \(\int_{\Gamma''\backslash\mathsf G_m}\sum_{0\ne y\in(c+aL)\gamma}f(y)\,d\mu''=a^{-Dm}\int f\) (A.4); (ii) Theorem R of R15 holds for every hereditary order without (h1\(^\ell\)); (iii) Theorem B holds for arbitrary Steinitz classes of \(M\) and \(L_0\), with the single constant \(\mathrm{covol}(ML_0)^{-1}\).
\item
  Not done: non-hereditary orders (the splittings B.1, B.3 may fail; one could try finite-index sublattices of product form and coset sums as in Part A, not attempted). Lattices \(\Lambda_W\otimes_{\mathcal O}L\) with \(L\) NON-free: the family \(\{L_0\gamma\}\), \(\Gamma=\mathrm{Stab}(L_0)\), is covered by Theorem B for the mean value side (\(\{X\in\Lambda_WL_0:\mathrm{col}(X)=U\}=\{B_UY:Y\in M_UL_0,\mathrm{rk}Y=j\}\) needs \((\Lambda_W\cap U)L_0=\Lambda_WL_0\cap U^{k}\), true when \(\Lambda_W\cap U\) is a direct summand of \(\Lambda_W\), which holds by (Her)); the covolumes then carry the factor \(\mathrm{covol}(M_UL_0)\) instead of \(c_U^k\)-type expressions and Lemma 1.4 would have to be redone. Not needed for Theorem E\textquotesingle.
\end{enumerate}
